\section{Causal dispersive material representation}
\label{sec:dispersive-materials}

Arbitrary media are useful only if their constitutive representation remains
causal, passive, and differentiable. A table of unrelated complex values at
each frequency is not an acceptable production material model.

\subsection{Frequency-domain constitutive law}

For an isotropic electric response,
\[
  D(\omega)=\epsilon_0\epsilon_r(\omega)E(\omega),
\]
and similarly for magnetic response. A causal rational representation is
written as
\[
 \epsilon_r(s)
 =
 \epsilon_\infty
 +\sum_{k=1}^{N_\epsilon}
   \frac{a_k}{s-p_k}
 +\sum_{k=1}^{N_\epsilon}
   \frac{\overline{a_k}}{s-\overline{p_k}},
 \qquad \Rea p_k<0 .
\]
Equivalent real Debye, Lorentz, or Drude parameterizations are allowed when they
match the material class. Conductivity may be represented separately or
absorbed into complex permittivity under a declared convention.

\subsection{Passivity}

A fitted material law is accepted only if the associated constitutive transfer
is passive over the declared band. Stable poles alone are insufficient.
Passivity is imposed either by a known passive pole/residue parameterization or
by a state-space realization whose storage function is positive.

For anisotropic material, the tensor transfer
\[
  \epsilon(s)=\epsilon(s)^\T
\]
for reciprocal media is parameterized through structured matrix residues or a
passive state realization; independently fitting each tensor entry is
forbidden if it can violate reciprocity or positivity.

\subsection{Auxiliary material states}

Dispersive media can be represented in time domain by local auxiliary states
$x_m$:
\[
 \dot x_m=A_m x_m+B_m E,\qquad
 P=C_m x_m+D_m E.
\]
This form couples naturally to broadband electromagnetic macromodels and avoids
reconstructing a frequency table at every time step.

\subsection{Temperature dependence}

Let $\vartheta$ denote temperature or a small vector of thermodynamic state
variables. Material parameters are smooth structured functions:
\[
  p_k=p_k(\vartheta),\qquad
  a_k=a_k(\vartheta),\qquad
  \epsilon_\infty=\epsilon_\infty(\vartheta).
\]
The interpolation is performed in stable/passive coordinates. Interpolating
complex permittivity entrywise between two temperatures is not sufficient if it
leaves the passive material manifold.

\subsection{Material API}

Every material definition declares:
\begin{enumerate}
\item the valid frequency and temperature domain;
\item scalar/tensor isotropy class and reciprocity;
\item causal/passive constitutive realization;
\item derivatives or differentiable evaluation with respect to frequency,
      temperature, and trainable material parameters;
\item uncertainty or data provenance when fitted from measurement.
\end{enumerate}

Material data outside the declared domain trigger extrapolation status and may
force CERTIFIED/REFERENCE evaluation.

\subsection{Background Green kernels}

A homogeneous or layered Green kernel uses the same causal constitutive law.
Kernel caching therefore depends on the material realization and frequency, not
on a material name. Two differently named materials with identical
constitutive realizations generate the same physics cache identity.
